The Routing Protocol for Low-Power and Lossy Networks (RPL) underpins large-scale Internet of Things (IoT) deployments, yet its open control plane exposes constrained nodes to internal attacks such as rank manipulation, selective forwarding, wormhole tunneling, and DAO flooding.
Existing intrusion detection systems (IDS) for RPL often depend on a centralized border router, generate excessive peer-to-peer alerts in flat distributed designs, or ignore residual energy and traffic criticality.
This paper introduces RPL-ClusterIDS, a clustering-based hierarchical distributed IDS that organizes detection around dynamically formed, energy-aware clusters shaped by normalized residual energy, topological stability, and traffic-priority class.
Ordinary members execute ultra-lightweight rules and report compact local anomaly scores, while cluster heads perform rule-based two-stage threshold verification under a context-adaptive policy that modulates cluster granularity and monitoring depth.
The system is implemented in Contiki-NG and evaluated through a Cooja pilot campaign (21 seeds for the CLUSTERIDS variant, 3 seeds for the ablation study; 5 attack scenarios on a 50-node grid).
In this pilot campaign, RPL-ClusterIDS reached 80.6\% campaign-level detection rate\footnote{Campaign-level detection rate is defined as at least one confirmed alert during an attack period.} (95.1\% on individual attack scenarios) with 0.50\% false-positive rate and 5--7\% CPU overhead, while the B1 baseline remained at 0\% detection rate. Per-interval detection rate averaged 1.1\% across modes (see Fig.~\ref{fig:mode-sensitivity}). These findings should be interpreted as promising preliminary evidence rather than definitive performance claims.
The full artifact package, including firmware, campaign scripts, parsed datasets, and analysis pipeline, is publicly available at: \url{https://github.com/madani-belacel/IDS-IOT}
